\section{Conclusion}
\label{sec:conclusion}

EconIRL provides twelve interchangeable estimators spanning structural econometrics and inverse reinforcement learning, unified by a common \texttt{Estimator} protocol and shared post-estimation infrastructure for standard errors, identification diagnostics, and counterfactual simulation. The tabular case study in Section~\ref{sec:experiments} confirms the theoretical equivalence of NFXP, CCP, and MCE-IRL, with all three recovering the same parameters and producing identical policy performance on the \citet{rust1987} bus engine problem. The scaling study shows that sieve and neural methods extend to high-dimensional state spaces without sacrificing parameter accuracy, maintaining cosine similarity above 0.99 to the true reward vector even as the number of states grows into the thousands. The nonlinear reward study demonstrates where deep IRL adds value over linear models, recovering reward curvature that linear specifications miss entirely.

Several limitations remain. All estimators except TD-CCP require known or consistently estimated transition matrices, which limits applicability in settings where the state transition process is complex or partially observed. The adversarial estimators, particularly AIRL, lack standard errors because the instability of GAN-style minimax training makes the asymptotic distribution of the discriminator parameters difficult to characterize. GLADIUS exhibits systematic bias in the IRL setting due to state-dependent identification constants that leak through asymmetric transitions, overestimating operating costs by approximately 40 percent on the bus engine benchmark regardless of network architecture or training duration.

Future extensions include dynamic games with strategic interaction \citep{aguirregabiriaMira2002}, unobserved heterogeneity via finite mixture and random coefficient models, and continuous-state deep estimators with valid inference guarantees. Integration with automatic differentiation through equilibrium constraints would enable efficient computation of standard errors for the adversarial estimators, addressing the most significant gap in the current inference pipeline. The package is MIT licensed and available at \url{https://github.com/econirl/econirl}.
